We built calibrated per-ball models of expected runs and wicket probability for
the T20 death overs on 223{,}678 deliveries and decomposed their predictive skill
into game state and a historical-rate encoding of player identity. We find no
evidence that this identity encoding improves on a state-only model out-of-sample:
the state-only model is the best-calibrated wicket
model, and adding batter/bowler identity worsens held-out log-loss and significantly
raises runs RMSE in all six rolling test seasons, even though SHAP attributes the
majority of explanatory mass to identity. A random-player placebo shows this
attribution is faithful but does not generalize --- a dissociation we trace to
overfitting of non-stationary form, with the caveat that it concerns our
historical-rate encoding of identity rather than identity in principle. Pressure
history violates the Markov assumption for scoring intensity (RMSE $+0.0065$,
dot-run sign negative) but not for wicket risk; this effect is associational and
plausibly confounded with bowler quality. Our headline is the momentum
measurement: importing the Miller--Sanjurjo finite-sequence bias correction to
per-ball cricket --- to our knowledge the first time the correction itself has been
applied at the delivery level --- we put the apparent-momentum illusion at
$+0.121$: the naive $-0.114$ estimate is almost exactly cancelled by a $-0.121$
finite-sequence bias, the corrected estimate is $+0.007$, and no cell survives a
state-stratified null. We report this as a measurement of the bias size --- the
quantity we are powered to estimate --- so the apparent ``cold hand'' is, to the
resolution we can measure, an artifact of the estimator rather than a property of
the data.

These results matter for both modeling and interpretation. They support
identity-agnostic, state-based per-ball models for the death overs as a calibrated
default, caution against reading SHAP attribution as out-of-sample importance, and
give the cricket pressure literature a bias correction it has lacked. Two concrete
next steps are a richer, interaction-aware identity representation --- to test
whether identity can be made to help rather than hurt --- and extending the
corrected momentum analysis to other phases and to women's T20 to map how the size
of the apparent-momentum illusion changes with the regime.
